% Options for packages loaded elsewhere
\PassOptionsToPackage{unicode}{hyperref}
\PassOptionsToPackage{hyphens}{url}
%
\documentclass[
  12pt,
]{article}
\usepackage{amsmath,amssymb}
\usepackage{iftex}
\ifPDFTeX
  \usepackage[T1]{fontenc}
  \usepackage[utf8]{inputenc}
  \usepackage{textcomp} % provide euro and other symbols
\else % if luatex or xetex
  \usepackage{unicode-math} % this also loads fontspec
  \defaultfontfeatures{Scale=MatchLowercase}
  \defaultfontfeatures[\rmfamily]{Ligatures=TeX,Scale=1}
\fi
\usepackage{lmodern}
\ifPDFTeX\else
  % xetex/luatex font selection
\fi
% Use upquote if available, for straight quotes in verbatim environments
\IfFileExists{upquote.sty}{\usepackage{upquote}}{}
\IfFileExists{microtype.sty}{% use microtype if available
  \usepackage[]{microtype}
  \UseMicrotypeSet[protrusion]{basicmath} % disable protrusion for tt fonts
}{}
\makeatletter
\@ifundefined{KOMAClassName}{% if non-KOMA class
  \IfFileExists{parskip.sty}{%
    \usepackage{parskip}
  }{% else
    \setlength{\parindent}{0pt}
    \setlength{\parskip}{6pt plus 2pt minus 1pt}}
}{% if KOMA class
  \KOMAoptions{parskip=half}}
\makeatother
\usepackage{xcolor}
\usepackage[margin=1in]{geometry}
\usepackage{graphicx}
\makeatletter
\def\maxwidth{\ifdim\Gin@nat@width>\linewidth\linewidth\else\Gin@nat@width\fi}
\def\maxheight{\ifdim\Gin@nat@height>\textheight\textheight\else\Gin@nat@height\fi}
\makeatother
% Scale images if necessary, so that they will not overflow the page
% margins by default, and it is still possible to overwrite the defaults
% using explicit options in \includegraphics[width, height, ...]{}
\setkeys{Gin}{width=\maxwidth,height=\maxheight,keepaspectratio}
% Set default figure placement to htbp
\makeatletter
\def\fps@figure{htbp}
\makeatother
\setlength{\emergencystretch}{3em} % prevent overfull lines
\providecommand{\tightlist}{%
  \setlength{\itemsep}{0pt}\setlength{\parskip}{0pt}}
\setcounter{secnumdepth}{5}
\usepackage{url}
\usepackage{setspace}
%\singlespacing
%\onehalfspacing
%\doublespacing
\usepackage{lineno}
%\linenumbers
\usepackage[belowskip=0pt,aboveskip=0pt]{caption}
\usepackage{relsize}
\newcommand{\Fmsy}{\ensuremath{F_{\text{MSY}}}\xspace}
\newcommand{\Fspr}[1]{\ensuremath{F_{\text{{#1}\%}}}\xspace}
\newcommand{\afrb}{Alaska Fishery Research Bulletin\xspace}
\newcommand{\ajms}{African Journal of Marine Science\xspace}
\newcommand{\amb}{Advances in Marine Biology\xspace}
\newcommand{\bms}{Bulletin of Marine Science\xspace}
\newcommand{\bjssf}{Bulletin of the Japanese Society of Scientific Fisheries\xspace}
\newcommand{\cb}{Conservation Biology\xspace}
\newcommand{\cjfas}{Canadian Journal of Fisheries and Aquatic Sciences\xspace}
\newcommand{\ea}{Ecological Applications\xspace}
\newcommand{\eer}{Evolutionary Ecology Research\xspace}
\newcommand{\elet}{Ecology Letters\xspace}
\newcommand{\emod}{Ecological Modelling\xspace}
\newcommand{\ebf}{Environmental Biology of Fishes\xspace}
\newcommand{\ff}{Fish and Fisheries\xspace}
\newcommand{\fo}{Fisheries Oceanography\xspace}
\newcommand{\fr}{Fisheries Research\xspace}
\newcommand{\fb}{Fishery Bulletin\xspace}
\newcommand{\ijms}{ICES Journal of Marine Science\xspace}
\newcommand{\iccat}{Collective Volume of Scientific Papers ICCAT\xspace}
\newcommand{\jae}{Journal of Animal Ecology\xspace}
\newcommand{\jai}{Journal of Applied Ichthyology\xspace}
\newcommand{\jdc}{Journal Du Conseil International Pour L'exploration De La Mer\xspace}
\newcommand{\jdcp}{Journal Du Conseil Permanent International Pour L'exploration De La Mer\xspace}
\newcommand{\jembe}{Journal of Experimental Marine Biology and Ecology\xspace}
\newcommand{\jfb}{Journal of Fish Biology\xspace}
\newcommand{\jsr}{Journal of Sea Research\xspace}
\newcommand{\jtb}{Journal of Theoretical Biology\xspace}
\newcommand{\jfrbc}{Journal of the Fisheries Research Board of Canada\xspace}
\newcommand{\jnwafs}{Journal of Northwest Atlantic Fisheries Science\xspace}
\newcommand{\mcf}{Marine and Coastal Fisheries: Dynamics, Management, and Ecosystem Science\xspace}
\newcommand{\mb}{Marine Biology\xspace}
\newcommand{\meps}{Marine Ecology Progress Series\xspace}
\newcommand{\mfr}{Marine Fisheries Review\xspace}
\newcommand{\mpb}{Marine Pollution Bulletin\xspace}
\newcommand{\najfm}{North American Journal of Fisheries Management\xspace}
\newcommand{\nzjmfr}{New Zealand Journal of Marine and Freshwater Research\xspace}
\newcommand{\pnas}{Proceedings of the National Academy of Sciences USA\xspace}
\newcommand{\rpvrciemm}{Rapports et Proc\`es-Verbaux des R\'eunions. Conseil Internationale pour l'Exploration de la Mer\xspace}
\newcommand{\rpvrcpiemm}{Rapports et Proc\`es-Verbaux des R\'eunions. Conseil Permanent Internationale pour l'Exploration de la Mer\xspace}
\newcommand{\rfbf}{Reviews in Fish Biology and Fisheries\xspace}
\newcommand{\sajms}{South African Journal of Marine Science\xspace}
\newcommand{\tafs}{Transactions of the American Fisheries Society\xspace}

\newcommand{\anzjs}{Australian \& New Zealand Journal of Statistics\xspace}
\newcommand{\as}{Applied Statistics\xspace}
\newcommand{\csda}{Computational Statistics \& Data Analysis\xspace}
\newcommand{\ees}{Environmental and Ecological Statistics\xspace}
\newcommand{\jas}{Journal of Applied Statistics\xspace}
\newcommand{\jabes}{Journal of Agricultural, Biological, and Environmental Statistics\xspace}
\newcommand{\jasa}{Journal of the American Statistical Association\xspace}
\newcommand{\jrssb}{Journal of the Royal Statistical Society. Series B\xspace}
\newcommand{\sm}{Statistics in Medicine}

\usepackage{float}
\usepackage{amsmath}
\usepackage{pdflscape}
\usepackage{hyperref}
\usepackage{xspace}
\usepackage{bm}
\usepackage{caption,graphics}
\usepackage{graphicx}
\usepackage{makecell}
\usepackage{lineno}
\linenumbers
\renewcommand\figurename{Fig.}
\renewcommand{\figureautorefname}{Fig.}
\captionsetup{labelsep=period, singlelinecheck=false}
\newcommand{\changesize}[1]{\fontsize{#1pt}{#1pt}\selectfont}
\renewcommand{\arraystretch}{1.5}
\renewcommand\theadfont{}
\usepackage{booktabs}
\usepackage{longtable}
\usepackage{array}
\usepackage{multirow}
\usepackage{wrapfig}
\usepackage{float}
\usepackage{colortbl}
\usepackage{pdflscape}
\usepackage{tabu}
\usepackage{threeparttable}
\usepackage{threeparttablex}
\usepackage[normalem]{ulem}
\usepackage{makecell}
\usepackage{xcolor}
\ifLuaTeX
  \usepackage{selnolig}  % disable illegal ligatures
\fi
\usepackage[]{natbib}
\bibliographystyle{cjfas2.bst}
\IfFileExists{bookmark.sty}{\usepackage{bookmark}}{\usepackage{hyperref}}
\IfFileExists{xurl.sty}{\usepackage{xurl}}{} % add URL line breaks if available
\urlstyle{same}
\hypersetup{
  pdftitle={Random effects on numbers-at-age transitions implicitly account for movement dynamics and improve stock assessment and management},
  pdfauthor={Chengxue Li\^{}\{1,6,7,*\}; Jonathan J. Deroba1; Aaron M. Berger2; Daniel R. Goethel3; Brian J. Langseth4; Amy M. Schueller5; Timothy J. Miller1},
  hidelinks,
  pdfcreator={LaTeX via pandoc}}

\title{Random effects on numbers-at-age transitions implicitly account
for movement dynamics and improve stock assessment and management}
\author{Chengxue Li\(^{1,6,7,*}\) \and Jonathan J.
Deroba\textsuperscript{1} \and Aaron M.
Berger\textsuperscript{2} \and Daniel R.
Goethel\textsuperscript{3} \and Brian J.
Langseth\textsuperscript{4} \and Amy M.
Schueller\textsuperscript{5} \and Timothy J. Miller\textsuperscript{1}}
\date{}

\begin{document}
\maketitle

\(^{1}\) NOAA, Northeast Fisheries Science Center, 166 Water Street,
Woods Hole, MA 02543, USA

\(^{2}\) NOAA, Northwest Fisheries Science Center, 2032 SE OSU Drive,
Newport, OR 97365, USA

\(^{3}\) NOAA, Alaska Fisheries Science Center, 17109 Point Lena Loop
Road, Juneau, AK 99801, USA

\(^{4}\) NOAA, Northwest Fisheries Science Center, 2725 Montlake Blvd E,
Seattle, WA 98112, USA

\(^{5}\) NOAA, Southeast Fisheries Science Center, 101 Pivers Island
Road, Beaufort, NC 28516, USA

\(^{6}\) School of Marine and Atmospheric Sciences, Stony Brook
University, Stony Brook, NY 11794, USA

\(^{7}\) Saltwater Inc., 733 N Street, Anchorage, AK 99501, USA

\(^{*}\) Corresponding author:
\href{mailto:chengxue.li@noaa.gov}{\nolinkurl{chengxue.li@noaa.gov}}

\vspace{0.5cm}

\textbf{ORCID:}

Chengxue Li:
\href{https://orcid.org/0000-0001-7518-7234}{0000-0001-7518-7234}

Aaron M. Berger:
\href{https://orcid.org/0000-0002-1408-7122}{0000-0002-1408-7122}

Daniel R. Goethel
\href{https://orcid.org/0000-0001-7121-6181}{0000-0001-7121-6181}

Brian J. Langseth
\href{https://orcid.org/0000-0002-9901-6146}{0000-0002-9901-6146}

Timothy J. Miller:
\href{https://orcid.org/0000-0003-1411-1206}{0000-0003-1411-1206}

\pagebreak

\hypertarget{abstract}{%
\section*{Abstract}\label{abstract}}
\addcontentsline{toc}{section}{Abstract}

Implementing operational assessment models that account for spatial
structure and movement dynamics is challenging, especially with limited
tagging data. Random effects on numbers-at-age (NAA) transitions in
state-space models offer a potential solution to circumvent direct
movement estimation by attributing movement variation to NAA random
effects. However, whether this approach reliably achieves desirable
management outcomes remains unclear. In this study, we conducted a
management strategy evaluation that emulated a generic medium-lived fish
that exhibit natal homing dynamics, using assessment models with varying
levels of spatial complexity. We compared the performance of each
spatial implementation with and without NAA random effects to evaluate
their effectiveness in achieving management outcomes. Our results showed
that models with NAA random effects consistently outperformed those
without, although the benefits of NAA random effects degraded at high
rates of movement. Therefore, NAA random effects could serve as a
practical intermediate solution when explicit movement modeling is not
feasible due to insufficient movement information. Our findings suggest
that incorporating NAA random effects should be a default starting point
in state-space stock assessments.

\textbf{Keywords}: state-space models, management strategy evaluation,
movement dynamics, numbers-at-age transitions, random effects

\pagebreak

\hypertarget{introduction}{%
\section{Introduction}\label{introduction}}

Random effects on numbers-at-age (NAA) transitions introduce flexibility
into a state-space stock assessment that can partially represent
misspecification of unmodeled processes, such as movement, natural
mortality, selectivity, or misreported catch
\citep{Nielsen2014, Cadigan2016, Albertsen2018, Perretti2020, Stock2021a, Stock2021b}.
Simulation-estimation studies have consistently concluded that
incorporating NAA random effects results in less biased estimates of
management quantities (e.g., biomass, reference points), even when NAA
random effects were included in the assessment model but absent in the
operating model that generated the data \citep{Li2024}. Previous studies
suggested that NAA deviations can be interpreted as migrations or
irregular natural mortality \citep{Frisk2010, Gudmundsson2012}. However,
the extent to which NAA random effects improve assessment results and
management outcomes when accounting for movement has not been explicitly
tested within the state-space modeling framework, despite claims that
they can do so \citep{Gudmundsson2012, Nielsen2014, Stock2021a}.

The benefits of accounting for spatial structure and movement dynamics
in stock assessments have become increasingly apparent
\citep{Goethel2021, Bosley2022, Goethel2023, Berger2024} and are
necessary as climate change continues to drive fish movement, alter fish
biology and ecology, and affect habitat availability and suitability
\citep{Portner2010, Ciannelli2013}. Acknowledging biocomplexity and
spatial heterogeneity in fish distribution and demographics challenges
the panmictic assumptions of traditional stock assessment models, which
treat the stock as as a single population that is uniformly distributed,
demographically homogeneous, closed to migration, and reproductively
isolated \citep{Beverton1957, Cadrin2020}. In practice, these models
represent heterogeneous stocks by estimating population totals based on
average weights, maturities, and mortality rates, thereby smoothing over
underlying spatial and demographic variation. Assessment models that do
not account for spatial structure and movement dynamics, when in fact
they are present, can produce biased estimates of population quantities,
potentially failing to achieve management objectives
\citep{Bosley2019, Goethel2023}. Including spatial structure and
movement dynamics in stock assessments, however, can be challenging
\citep{Punt2019, Punt2020, Goethel2022, Berger2021, Berger2024}, and in
some cases may not be practical due to jurisdictional or management
constraints \citep{Kerr2017}, low quality and quantity of
geo-partitioned data \citep{Cadrin2020, Punt2020}, inadequate tagging
information \citep{Goethel2011, Goethel2019}, or limited knowledge about
population origin and connectivity \citep{Berger2017}.

The ``fleets-as-areas'' approach is often used as a practical starting
point in spatial assessments
\citep{Berger2012, Waterhouse2014, Lee2017, Bosley2022}. This approach
assumes the population is uniformly distributed across its range; any
apparent spatial heterogeneity is represented indirectly by defining
area-specific fleets, with differences in age compositions attributed to
fishing selectivity rather than explicit spatial structure. However,
differences in estimated selectivity across areas may not reflect actual
differences in the underlying gear- or fleet-specific selectivity
patterns. Instead, such differences may arise from spatial heterogeneity
in the population itself (e.g., age structure differences caused by
movement, source--sink dynamics, or environmentally driven shifts in
distribution) \citep{Cope2011, Berger2012, Hurtado2014, O'Boyle2016}.
When such dynamics are present, the implicit mismatch between the
assumption (spatially varying selectivity) and the underlying process
(spatially varying distribution) can undermine the performance of the
assessment--management framework \citep{Punt2016, Punt2017}.

Explicitly specifying population structure and spatial units, but
ignoring movement, provides an alternative approach for improving
population estimates and derived management parameters in a spatial
context, particularly when adult movement is infeasible to estimate
\citep{Cadrin2019, Goethel2023}. Evidence suggests that incorporating
spatially referenced parameters within an assessment model to account
for the net effects of movement on observations of population structure
at the local scale is both practical and broadly applicable to many
fisheries \citep{Berger2012, Goethel2023}. This approach is particularly
useful in cases where species biology varies geographically due to local
biophysical conditions, habitat availability and suitability, or harvest
patterns \citep{Goethel2023}, and when adult movement is limited, with
larval dispersal driving spatial structure \citep{Kerr2017}.

Incorporating NAA random effects in state-space models represents a step
forward in model complexity, where spatial structure is implicitly or
explicitly defined, and movement is handled implicitly through NAA
random effects. Evaluating whether NAA random effects can improve
management outcomes in stock assessments that exclude explicit movement
modeling is crucial. Using NAA random effects as an ``add-on'' proxy for
movement simplifies movement dynamics and reduces the number of
parameters estimated in the assessment model. However, despite
simplifying parameter estimation, NAA random effects do not
mechanistically capture the underlying movement dynamics of fish
populations. Instead, movement variability is absorbed into the NAA
random effects structure, which can lead to misinterpretations of
population processes, model misspecification, and biased estimates of
stock status and productivity.

Thus, NAA random effects may provide a more tractable solution, but the
trade-off between model complexity and biological realism must be
carefully considered, especially in the context of management
performance. One effective tool to evaluate whether management
procedures achieve management objectives (e.g., sustainable catch) is
through management strategy evaluation (MSE)
\citep{Punt2017,Keymer2000}. The MSE approach in a spatial context
involves considering various candidate spatial population structures,
data collection methods, estimation model structures, harvest control
rules, and their interactions and feedbacks
\citep{Kerr2014, Goethel2016}.

In this study, we use a MSE with closed-loop simulations to evaluate the
performance of several spatial assessment model configurations that do
not explicitly account for movement but include NAA random effects,
while the operating model defining the true population dynamics includes
movement. We examine the utility of incorporating NAA random effects
across increasing levels of assessment model complexity (panmictic,
spatially implicit, and spatially explicit) and assess their influence
on estimated assessment quantities and management outcomes under varying
movement dynamics and historical fishing pressure.

\hypertarget{methods}{%
\section{Methods}\label{methods}}

\hypertarget{model-overview}{%
\subsection{Model Overview}\label{model-overview}}

The Woods Hole Assessment Model (WHAM,
\url{https://timjmiller.github.io/wham}), a state-space assessment
model, was used as the basis of the operating model (OM) and estimation
model (or assessment model; EM) in the MSE framework. WHAM can
incorporate time- and/or age-varying process errors on population and
fishery processes, including recruitment, NAA, natural mortality,
movement, fishing selectivity, and survey catchability, and it has
ability to link environmental covariates to specific processes
\citep{Stock2021a}. More recently, WHAM has been extended to incorporate
spatial dynamics, allowing populations to be modeled with spatial
heterogeneity in biology, movement, and fleet dynamics
\citep{Miller2025review}. Here we give a minimal overview necessary for
the present study (see Methods in the supplementary file for details).
We used the SPASAM-MSE package
(\url{https://lichengxue.github.io/whamMSE/}), a comprehensive MSE
toolbox developed upon WHAM \citep{Liinreview2}, to conduct MSE and
explore potential implications of management strategies.

\hypertarget{operating-models-oms}{%
\subsection{Operating Models (OMs)}\label{operating-models-oms}}

Two operating models (OMs) were developed that represent the true
dynamics of the system. Each was conditioned using parameters for fish
with a medium-lived (maximum age = 12) life history (Wiedenmann et al.,
2015) and differed only in the rates of movement assumed (high and low)
between populations and regions (\autoref{para_tab}). The spatial domain
included two regions, each of which is the natal home of one population
but can also be occupied by the other population outside the spawning
season. Population and fishery parameters were consistent across the
domain, with exceptions for stock-recruitment relationship and fishing
selectivity (\autoref{para_tab}).

\hypertarget{fishing-history}{%
\subsubsection{Fishing History}\label{fishing-history}}

For each OM, three alternative fishing histories (see
\autoref{F_scenario_table}) were used during the simulation historical
period: 1) \(\text{F}_\text{A}\): both regions experience overfishing
(2.5\(F_{MSY}\)) for the first 15 years and fishing at \(F_{MSY}\) for
the last 15 years; 2) \(\text{F}_\text{U}\): both regions experience
constant overfishing (2.5\(F_{MSY}\)) for the entire 30 years; and 3)
\(\text{F}_\text{C}\): region 1 experiences overfishing (2.5\(F_{MSY}\))
for the first 15 years and fishing at \(F_{MSY}\) for the last 15 years,
while region 2 follows the opposite fishing history.

\hypertarget{biological-information}{%
\subsubsection{Biological Information}\label{biological-information}}

The Beverton-Holt stock-recruitment relationship was assumed for both
populations, with slightly different productivity (\({\alpha}\))
parameters to create some discrepancy in biology between the two
populations. Weight-at-age, maturity-at-age, and natural mortality were
assumed to be the same between the two populations. Initial
numbers-at-age for each population were specified at the equilibrium
distribution fishing at \({F}\) = 0.1 in each region with an initial
recruitment of \(e^{10}\).

\hypertarget{numbers-at-age-naa-transitions}{%
\subsubsection{Numbers-at-age (NAA)
Transitions}\label{numbers-at-age-naa-transitions}}

For simplicity, we present the one-population, one-region version of the
NAA transition equations. When the number of populations (\(n_s > 1\))
or regions (\(n_r > 1\)) increases, the formulation extends to higher
dimensions by incorporating population- and region-specific random
effects \citep{Miller2025review}. The numbers-at-age transitions are:

\begin{equation}
\log(N_{a,y}) =
\begin{cases}
\log(f(\text{SSB}_{y-1})) + \epsilon_{1,y}, & a = 1 \\[-0.3em]
& \epsilon_{1,y} \sim N(0,\sigma_{Rec}^2) \\
\log(N_{a-1,y-1} e^{-Z_{a-1,y-1}}) + \epsilon_{a,y}, & 1 < a < A \\[-0.3em]
& \epsilon_{a,y} \sim N(0,\sigma_{NAA}^2) \\
\log(N_{A-1,y-1} e^{-Z_{A-1,y-1}} + N_{A,y-1} e^{-Z_{A,y-1}}) + \epsilon_{A,y}, & a = A \\[-0.3em]
& \epsilon_{A,y} \sim N(0,\sigma_{NAA}^2)
\end{cases}
\end{equation}

where \(N_{a,y}\) is the abundance at age \(a\) in year \(y\),
\(Z_{a,y} = F_{a,y} + M_{a,y}\) is total mortality, \(f(\cdot)\) is the
stock--recruitment function, \(A\) defines the plus-group, and
\(\epsilon\) are lognormal process errors. In this study, recruitment
and NAA transitions in the OM were treated as independent and
identically distributed (IID) random effects, with
\(\sigma_{Rec} = 0.5\) for recruitment and \(\sigma_{NAA} = 0.2\) for
age transitions across populations and regions.

\hypertarget{movement-dynamics}{%
\subsubsection{Movement Dynamics}\label{movement-dynamics}}

Four seasons (spring, summer, autumn, and winter) were included in the
OMs to represent seasonal migrations characteristic of natal homing
dynamics. Specifically, the two reproductively isolated populations were
assumed to return to their natal regions at the end of summer, spawn in
autumn, and migrate to other regions during spring and winter (e.g., for
feeding migrations). Two different movement scenarios (low and high)
were included in this study (see \autoref{fig:movement_dynamics}). The
high movement scenario (\(\text{OM}_{\text{high}}\)) allows both
populations to move between regions outside the spawning season. The
movement rate from region 1 to region 2 was assumed to be much higher
than the rate in the opposite direction, reflecting a degree of
source-sink dynamics, that could reflect habitat suitability
\citep{Keymer2000}. The low movement scenario
(\(\text{OM}_{\text{low}}\)) allows population 1 to move from region 1
to region 2 and back from region 2 to region 1, whereas population 2
cannot move between regions. This type of movement represents strict
source-sink movement dynamics typical of some species (e.g., black sea
bass, \emph{Centropristis striata}, Miller et al., In Press). The
movement rate outside the spawning season was assumed to be time-varying
and treated as a random effect, with a standard deviation of 0.5 and an
AR1 autocorrelation coefficient of 0.5. In each simulation replicate, a
new realization of these movement rates was generated; examples are
shown in Figs. S1 and S16.

\hypertarget{fleet-and-survey}{%
\subsubsection{Fleet and Survey}\label{fleet-and-survey}}

A fishing fleet was assumed to operate in each region, with logistic
selectivity patterns differing slightly between the two fleets to
reflect spatial heterogeneity (\autoref{para_tab}). A survey was assumed
to be conducted in each region during the spawning season to provide the
most accurate index of \(SSB\) for each population. Both surveys share
the same selectivity curve, with a lower \(a_{50}\) than that of the
fleet, ensuring that adequate information on young fish can be
collected. A multinomial distribution with an effective sample size of
100 was assumed for the fleet and survey composition data, and a
log-normal distribution with a coefficient of variation (CV) of 0.1 was
used for aggregate catch and survey indices. We also conducted a
sensitivity analysis by increasing the survey CV to 0.3. Weight-at-age
was assumed to be the same across populations, regions, fleets, and
surveys.

\hypertarget{estimation-models-ems}{%
\subsection{Estimation Models (EMs)}\label{estimation-models-ems}}

EMs with varying population structures, complexities, and the inclusion
or exclusion of NAA random effects were explored
(\autoref{fig:EM_structure}; \autoref{EM_tabel}). For simplicity, EMs
that included NAA random effects are hereafter referred to as
\(\text{EMs}_{\scriptscriptstyle \text{NAA}}\), while EMs without NAA
random effects are referred to as
\(\text{EMs}_{\scriptscriptstyle \text{noNAA}}\).

Of the 10 EMs, the first 8 were paired (with and without NAA random
effects) and included a single spatially aggregated panmictic assessment
model (\(\text{PAN}_{\scriptscriptstyle \text{NAA}}\) and
\(\text{PAN}_{\scriptscriptstyle \text{noNAA}}\)), a single spatially
implicit assessment model
(\(\text{FAA}_{\scriptscriptstyle \text{NAA}}\) and
\(\text{FAA}_{\scriptscriptstyle \text{noNAA}}\)), two separate
individual assessment models by region
(\(\text{SEP}_{\scriptscriptstyle \text{NAA}}\) and
\(\text{SEP}_{\scriptscriptstyle \text{noNAA}}\)), and a spatially
disaggregated assessment model with no movement
(\(\text{SpD}_{\scriptscriptstyle \text{NAA}}\) and
\(\text{SpD}_{\scriptscriptstyle \text{noNAA}}\)). The difference
between SEP and SpD EMs is that SpD EMs use global SPR-based reference
points weighted by regional recruitment, while SEP EMs separately
calculate regional SPR-based reference points. In addition, SpD EMs
allow certain parameters (e.g., NAA, movement) to be shared across
regions, whereas SEP EMs require separate parameter sets for each
region. The remaining two EMs were spatially explicit with movement
either fixed (\(\text{SpE}_{\scriptscriptstyle \text{NAA,Fix}}\)) at the
true value or with mean movement estimated
(\(\text{SpE}_{\scriptscriptstyle \text{NAA,Est}}\)) using a prior
distribution (with mean = \(\mu_{True}\) and standard deviation = 0.5).
When NAA random effects were included in the EM, they were modeled as
independent and identically distributed (IID) normal (\(\mathcal{N}\))
random effects in log space, corresponding to the \(\epsilon\) terms in
Equation 1. In this case, only a single additional variance parameter
(\(\sigma_{NAA}\)) is estimated, representing variability in the NAA
transitions for that population within each region. Recruitment in all
EMs was likewise modeled as IID normal random effects in log space.

The EM with fixed movement
(\(\text{SpE}_{\scriptscriptstyle \text{NAA,Fix}}\)) was considered as
the best possible EM and was treated as the baseline for comparison.

\hypertarget{management-strategy-evaluation-mse}{%
\subsection{Management Strategy Evaluation
(MSE)}\label{management-strategy-evaluation-mse}}

A MSE was conducted to evaluate the performance of the EMs. MSEs are
typically implemented as closed-loop simulations, where population
dynamics, data collection, stock assessment, and management actions are
linked together. Management decisions feed back into the system and
affect future population trajectories, allowing the long-term
consequences of alternative strategies to be evaluated in a realistic,
iterative framework.

In this study, the closed-loop simulation included a 30-year historical
period followed by a 30-year feedback period. Preliminary results
indicated that a 30-year feedback period was sufficient to capture
differences in management outcomes among the EMs compared to a longer
feedback period. Assessments were conducted every 3 years, with each
assessment using the most recent 20 years of data available at that
time. Sensitivity analysis indicated small differences in convergence
rates and parameter estimates of the EMs when the number of years of
input data was increased beyond 20 \citep{Li2018}.

In the interim years between assessments, the population was projected
in the EM using the average selectivity-at-age, maturity-at-age,
weight-at-age, and natural mortality from the last 5 model years of the
most recent assessment, as well as the average recruitment across all
model years. For EMs with processes (e.g., recruitment, NAA, movement)
treated as random effects, only recruitment random effects were assumed
to continue while other process random effects were set to zero during
the projection years \citep{Stock2021a}. To avoid confounding effects
from differences in model specifications of process random effects and
spatial dynamics, we elected to allow only recruitment random effects to
persist into the projection period. This ensured consistency across EMs,
since all models included recruitment random effects.

Target catch in those interim years, was based on a constant fishing
mortality harvest control rule at \(F_{40\%}\). The SEP EM used
region-specific SPR-based reference points, whereas all other EMs used
global SPR-based reference points. For SpD and SpE EMs, additional
weights based on the mean recruitment of each population estimated by
the EM were applied when calculating the global SPR-based reference
points. In contrast, no such weights were used for PAN and FAA EMs, as
they assumed a single population and relied on a total estimated
recruitment. For the PAN EM, the total catch target was allocated to
region-specific catches using weights derived from the relative biomass
catch observed in each survey during the terminal year of the assessment
\citep{Bosley2019}. All other EMs provided region-specific catch advice
that was inherently based on regional \(F\)-at-age matrices, all of
which collectively contribute to the calculation of the global \(F\).
See Table S1 for more details about catch apportionment.

The annual catch advice from each EM was sequentially fed into the OMs
to update \(F\) and population dynamics. This closed-loop simulation was
performed iteratively until the end of the feedback period. The EMs were
first fit at the end of the 30-year historical period (year 30) and then
refit every three years during the 30-year feedback period, resulting in
10 EM fits per OM iteration. The final EM fit occurred in year 57, and
its catch advice was used to update the OM through the end of year 60.
If any of the 10 EMs failed to converge during a replicate, the results
for that entire replicate, across all EMs, were discarded. A new random
seed was selected, and the replicate was attempted again until 100
replicates successfully converged for all 10 EMs in the study. In total,
the MSE evaluated 60 scenarios comprising two states of nature (high or
low movement), three fishing histories, and 10 EMs.

\hypertarget{em-diagnostics}{%
\subsubsection{EM Diagnostics}\label{em-diagnostics}}

Performance was summarized over 100 converged runs to ensure equal
sample sizes across scenarios. Convergence rates for each EM were
calculated based on the first 100 realizations (including replicates
that failed to converge) for each OM and fishing history. In addition to
convergence rates, all 10 EMs were evaluated by comparing their
estimates of mean recruitment, recruitment standard deviation, and NAA
standard deviation, using all 1,000 converged models (100 replicates per
EM) for each OM and fishing history.

\hypertarget{simulation-estimation-results}{%
\subsubsection{Simulation-Estimation
Results}\label{simulation-estimation-results}}

To evaluate estimation bias of each EM, the relative bias of estimated
annual recruitments and \(SSB\) was calculated for each EM compared to
the true values from the simulated pseudo-dataset \(i\). The median
relative bias for each EM from the first assessment was calculated as:

\begin{equation}
\text{Relative Error}_{i} = \text{Median} \left( \frac{\hat{\theta}_{i,y}}{\theta_{i,y}} - 1 \right)
\end{equation}

where \(\theta(i,y)\) represents the true value for year \(y\) from the
simulated pseudo-dataset \(i\), and \(\hat{\theta}(i,y)\) is the
estimated value obtained from the EM fitted to the pseudo-dataset.

To compare parameter estimates associated with recruitment and NAA
processes, estimates of mean recruitment (\(\mu_{Rec}\)), recruitment
standard deviation (\(\sigma_{Rec}\)), and NAA standard deviation
(\(\sigma_{NAA}\)) were also collected from the EMs at the last
iteration of the closed-loop.

\hypertarget{mse-performance-metrics}{%
\subsubsection{MSE Performance Metrics}\label{mse-performance-metrics}}

Short-term (first 5 years of the feedback period) and long-term (last 5
years of the feedback period) performance metrics at both global and
regional scales were evaluated to quantify the performance of each EM.
These performance metrics include:

\begin{itemize}
\item
  Average catch: \(\frac{1}{n} \sum_{y=1}^{n} {\text{Catch}}_y\);
\item
  Average fully selected \({F}\): \(\frac{1}{n} \sum_{y=1}^{n} {F}_y\);
\item
  Average \({SSB}\): \(\frac{1}{n} \sum_{y=1}^{n} {SSB}_y\);
\end{itemize}

Here, catch refers to the catch advice projected by each EM. This
projected catch is treated as the realized catch in the OM, with no
observation error added (i.e., projected catch from the EM = actual
catch in the OM). When updating \(F\) in the OM, the true catch is used
without observation error, and the resulting \(SSB\) represents the true
\(SSB\) from the OM given that \(F\). Although observed catch values
with error are generated in the observation model and can be used for
comparison, they are not used in the calculation of the performance
metrics presented here.

In our study, \(\text{SpE}_{\scriptscriptstyle \text{NAA,Fix}}\) was
used as the baseline EM because it represents performance with correctly
specified movement, thereby reflecting the best possible performance.
For replicate \(i\), the relative difference between the performance
metrics of the baseline EM and other EMs was calculated as:

\begin{equation}
\text{Relative Difference }(i) = \frac{\text{Metric}_{EM_{x}}(i)}{\text{Metric}_{\text{SpE}_{\scriptscriptstyle \text{NAA,Fix}}}(i)} - 1
\end{equation}

This approach standardizes differences across performance metrics,
providing better visualization and comparison on the same scale.

Following performance metrics were also collected:

\begin{itemize}
\item
  Probability of \({F} > {F_{MSY}}\):
  \(\text{P}_{\text{F>F}_{\text{MSY}}} = \frac{1}{n} \sum_{y=1}^{n} I\left( \frac{F_y}{F_{MSY_{y}}} > 1 \right)\);
\item
  Probability of \({SSB} < {SSB_{MSY}}\):
  \(\text{P}_{\text{SSB<SSB}_{\text{MSY}}} = \frac{1}{n} \sum_{y=1}^{n} I\left( \frac{SSB_y}{SSB_{MSY_{y}}} < 1 \right)\);
\item
  Overfishing status in the terminal year of the feedback period:
  \(\text{OFG}_\text{terminal} = \frac{F_T}{F_{MSY_{T}}}\);
\item
  Overfished status in the terminal year of the feedback period:
  \(\text{OFD}_\text{terminal} = \frac{SSB_T}{SSB_{MSY_{T}}}\);
\item
  Average annual catch variation:
  \(\text{AACV} = \frac{\sum_{y=1}^{n} \left|\text{Catch}_{y} - \text{Catch}_{y-1}\right|}{\sum_{y=1}^{n} \text{Catch}_{y-1}}\).
\end{itemize}

The annual global reference points (\(F_{MSY_{y}}\) and
\(SSB_{MSY_{y}}\)) were used for calculating some of the above
performance metrics \citep{Miller2025review}, with global values
obtained as a weighted sum, where weights were based on the mean of
realized annual recruitments for each population within each replicate.

All performance metrics for the 10 EMs at the global scale were
collected to provide a holistic view of model performance. Metrics such
as catch, \(F\), \(SSB\), terminal-year \(SSB/SSB_{MSY}\) for each OM
and fishing history were normalized to scores between 0 and 1 using:

\begin{equation}
\text{Score } (i) = \frac{\text{Metric}_{EM_x}(i) - \min(\text{Metric}_{EM_{A}}(i))}{\max(\text{Metric}_{EM_{A}}(i)) - \min(\text{Metric}_{EM_{A}}(i))}, \text{with } A = \{1, 2, \dots, 10\}
\end{equation}

For metrics where higher values indicated poorer performance (i.e.,
\(\text{P}_{\text{F>F}_{\text{MSY}}}\),
\(\text{P}_{\text{SSB<SSB}_{\text{MSY}}}\),
\(\text{OFG}_\text{terminal}\), and AACV), normalization was adjusted to
ensure higher values consistently indicate better performance, as
follows:

\begin{equation}
\text{Score } (i) = 1- \frac{\text{Metric}_{EM_x}(i) - \min(\text{Metric}_{EM_{A}}(i))}{\max(\text{Metric}_{EM_{A}}(i)) - \min(\text{Metric}_{EM_{A}}(i))}, \text{with } A = \{1, 2, \dots, 10\}
\end{equation}

The median, 25th quantile, and 75th quantile of the above performance
metrics and scores were calculated and presented in the results.

\hypertarget{results}{%
\section{Results}\label{results}}

Our simulation--estimation results indicated the
\(\text{SpE}_{\scriptscriptstyle \text{NAA,Est}}\) EM and the
\(\text{SpE}_{\scriptscriptstyle \text{NAA,Fix}}\) EM (baseline EM)
performed similarly, both producing accurate estimates of recruitment
(Figs. S2 and S17) and \(SSB\) (Figs. S3 and S18), as well as comparable
estimates of model parameters (Figs. S4 and S19). In addition, our MSE
results showed that \(\text{SpE}_{\scriptscriptstyle \text{NAA,Est}}\)
performed similarly to the baseline EM for most performance metrics and
performed relatively better than the other EMs.

To simplify the results we focused on a subset of representative cases.
Specifically, we excluded further discussion of the
\(\text{SpE}_{\scriptscriptstyle \text{NAA,Est}}\) EM since its results
closely mirrored those of the baseline EM. We also emphasized results
for \(\text{OM}_{\text{high}}\) rather than \(\text{OM}_{\text{low}}\),
because both produced qualitatively consistent outcomes but
\(\text{OM}_{\text{high}}\) highlighted differences among EMs more
clearly. Results of \(\text{OM}_{\text{low}}\) generally aligned with
those of \(\text{OM}_{\text{high}}\), but performance metrics among EMs
were less distinct due to lower movement rates. Additionally,
\(\text{EMs}_{\scriptscriptstyle \text{NAA}}\) under the
\(\text{OM}_{\text{low}}\) scenario generally performed similarly to the
baseline EM. A brief summary of the \(\text{OM}_{\text{low}}\) results
is included at the end of this section, with further detail in the
supplementary material (Figs. S16--S35).

A sensitivity analysis was conducted by increasing the survey CV from
0.1 to 0.3. We found that the results remained broadly similar, likely
because the CV for aggregate catch was already relatively low (0.1), so
increasing the CV for survey data had little influence on parameter
estimates, as the model placed more weight on the catch data.
Importantly, since our study focuses on comparing estimation models with
and without NAA random effects, changing the survey CV does not affect
the main conclusion regarding whether including NAA random effects helps
to alleviate issues related to ignoring movement dynamics, as long as
data quality is applied consistently across all scenarios.

Throughout the scenarios, convergence rate was high (85--100\%) for all
cases.

\hypertarget{simulation-estimation-results-1}{%
\subsection{Simulation-estimation
Results}\label{simulation-estimation-results-1}}

Total recruitment estimate from PAN and FAA EMs were similar to the sum
of region-specific recruitment from SEP and SpD EMs, all of which were
slightly overestimated compared to the total recruitment estimated from
the baseline EM (Fig. S2). The mean recruitment parameter
(\(\mu_{Rec}\)), recruitment standard deviation (\(\sigma_{Rec}\)) and
NAA standard deviation (\(\sigma_{NAA}\)) from EMs without movement
included, appeared to compensate to account for the absence of movement
dynamics (Fig. S4). This pattern indicates that variation arising from
movement dynamics was instead absorbed by recruitment and NAA random
effects, effectively compensating for the absence of explicit movement
processes.

\hypertarget{mse-performance-metrics-1}{%
\subsection{MSE Performance Metrics}\label{mse-performance-metrics-1}}

\hypertarget{relative-difference-in-annual-catch-and-ssb}{%
\subsubsection{\texorpdfstring{Relative Difference in Annual Catch and
\(SSB\)}{Relative Difference in Annual Catch and SSB}}\label{relative-difference-in-annual-catch-and-ssb}}

Relative difference in annual catch over the 30-year feedback period
revealed larger discrepancies among EMs during the first 10 years
compared to later years. As time progressed,
\(\text{EMs}_{\scriptscriptstyle \text{NAA}}\) performed more similarly
to the baseline EM (Fig. S5). This trend was more pronounced at the
regional scale (Figs. S6-S7). Relative differences in annual \(SSB\) at
both global (Fig. S8) and regional scales (Figs. S9-S10) tended to
diverge from zero over time.
\(\text{EMs}_{\scriptscriptstyle \text{NAA}}\) maintained higher \(SSB\)
that was more similar to the baseline EM than
\(\text{EMs}_{\scriptscriptstyle \text{noNAA}}\).

\hypertarget{short-term-catch-f-and-ssb}{%
\subsubsection{\texorpdfstring{Short-term Catch, \(F\), and
\(SSB\)}{Short-term Catch, F, and SSB}}\label{short-term-catch-f-and-ssb}}

In the first five years of the feedback period,
\(\text{EMs}_{\scriptscriptstyle \text{NAA}}\) generally produced catch,
\(F\), and \(SSB\) more similar with smaller variations (i.e.,
inter-quantile ranges) to the baseline EM than
\(\text{EMs}_{\scriptscriptstyle \text{noNAA}}\) at both regional and
global scales (\autoref{fig:first_5yr_results}). SEP and SpD EMs with
NAA random effects generally produced results more similar to the
baseline EM than other EMs, and this was especially true for the
\(\text{SpD}_{\scriptscriptstyle \text{NAA}}\) EM
(\autoref{fig:first_5yr_results}).

Performance in catch, \(F\), and \(SSB\) also exhibited difference
across EMs with difference in model complexity. For example, PAN EMs
generally produced higher catch (with higher \(F\)) in region 1 and
lower catch (with lower \(F\)) in region 2, and as a result lower
\(SSB\) in region 1 and higher \(SSB\) in region 2, although the global
performance metrics were similar to the baseline EM
(\autoref{fig:first_5yr_results}). Despite high contrast in \(F\)
between regions when using PAN EMs, \(SSB\) showed only marginal
difference compared to the baseline (5\% less) in the short term
(\autoref{fig:first_5yr_results}).

\hypertarget{long-term-catch-f-and-ssb}{%
\subsubsection{\texorpdfstring{Long-term Catch, \(F\), and
\(SSB\)}{Long-term Catch, F, and SSB}}\label{long-term-catch-f-and-ssb}}

In the last five years of the feedback period, there was marginal
difference in median catch (despite variations being higher for
\(\text{EMs}_{\scriptscriptstyle \text{noNAA}}\)) at both regional and
global scales between \(\text{EMs}_{\scriptscriptstyle \text{NAA}}\) and
\(\text{EMs}_{\scriptscriptstyle \text{noNAA}}\) across differing levels
of model complexity (\autoref{fig:last_5yr_results}). However,
\(\text{EMs}_{\scriptscriptstyle \text{noNAA}}\) often produced
significantly high \(F\), particularly at the regional scale
(\autoref{fig:last_5yr_results}). The consistently high \(F\) in
\(\text{EMs}_{\scriptscriptstyle \text{noNAA}}\) resulted in lower
\(SSB\) compared to the baseline EM at both regional and global scales
(\autoref{fig:last_5yr_results}).

Spatially implicit and explicit EMs (i.e., FAA, SEP, and SpD EMs)
performed similarly, particularly at the global scale, all of which
produced \textasciitilde10\% lower \(SSB\) when NAA random effects were
included and \textasciitilde20\% lower \(SSB\) when NAA random effects
were not included, compared to the baseline EM
(\autoref{fig:last_5yr_results}). The simplest PAN EMs performed
slightly better than other EMs at the global scale, particularly under
the unsustainable and contrasting fishing histories (5-10\% lower global
\(SSB\)) (\autoref{fig:last_5yr_results}). However, the corresponding
\(F\) for region 1 produced by PAN EMs was 50\%-75\% higher than the
\(F\) in the baseline EM (\autoref{fig:last_5yr_results}).

\hypertarget{average-annual-catch-variation-aacv}{%
\subsubsection{Average Annual Catch Variation
(AACV)}\label{average-annual-catch-variation-aacv}}

Including NAA random effects frequently lowered AACV, particularly at
the regional scale, regardless of model complexity and fishing histories
(\autoref{fig:aacv}). Although AACV at the global scale was similar
across EMs for each fishing history, some benefits of including NAA
random effects in reducing AACV were still observed under the adaptive
fishing history (\autoref{fig:aacv}). AACV at the global scale was
particularly low when the \(\text{PAN}_{\scriptscriptstyle \text{NAA}}\)
EM was used (\autoref{fig:aacv}). At the regional scale, relatively low
AACV occurred when using FAA EMs, which even outperformed the baseline
EM, regardless of fishing histories (\autoref{fig:aacv}). In contrast,
PAN and SEP EMs generally exhibited high AACV at the regional scale
(\autoref{fig:aacv}).

\hypertarget{probability-of-overfished-and-overfishing}{%
\subsubsection{Probability of Overfished and
Overfishing}\label{probability-of-overfished-and-overfishing}}

In most cases, \(\text{EMs}_{\scriptscriptstyle \text{NAA}}\)
outperformed \(\text{EMs}_{\scriptscriptstyle \text{noNAA}}\) in
preventing populations from becoming overfished and overfishing at the
global scale (\autoref{fig:Overfished_OM1}).
\(\text{EMs}_{\scriptscriptstyle \text{NAA}}\) exhibited similar
performance in terms of the probability of \({SSB} < {SSB_{MSY}}\)
compared to the baseline EM, but performed worse in preventing
overfishing (\autoref{fig:Overfished_OM1}). The probability of
overfishing was lower when using the
\(\text{PAN}_{\scriptscriptstyle \text{NAA}}\) EM compared to other
\(\text{EMs}_{\scriptscriptstyle \text{NAA}}\), except for SpE EMs, as
indicated by the consistently smaller overfishing probabilities in
\autoref{fig:Overfished_OM1}.

\hypertarget{terminal-year-status}{%
\subsubsection{Terminal-year Status}\label{terminal-year-status}}

Populations were not overfished in the terminal year at the global
scale, regardless of NAA random effects, model complexity, and fishing
histories (Fig. S12). \(\text{EMs}_{\scriptscriptstyle \text{NAA}}\)
performed better than \(\text{EMs}_{\scriptscriptstyle \text{noNAA}}\)
in preventing global overfishing in the terminal year (Fig. S12).
Notably, none of \(\text{EMs}_{\scriptscriptstyle \text{noNAA}}\) could
prevent global overfishing (Fig. S12).

\hypertarget{holistic-view-of-performance}{%
\subsubsection{Holistic View of
Performance}\label{holistic-view-of-performance}}

In general, \(\text{EMs}_{\scriptscriptstyle \text{NAA}}\) outperformed
\(\text{EMs}_{\scriptscriptstyle \text{noNAA}}\) across most metrics,
except for short-term catch (\autoref{fig:holistic_OM1}).
\(\text{SEP}_{\scriptscriptstyle \text{noNAA}}\) produced the highest
global catch under adaptive and unsustainable fishing histories, while
\(\text{PAN}_{\scriptscriptstyle \text{noNAA}}\) produced the highest
global catch under the contrasting fishing history
(\autoref{fig:holistic_OM1}). Relatively complex
\(\text{EMs}_{\scriptscriptstyle \text{noNAA}}\), such as
\(\text{SEP}_{\scriptscriptstyle \text{noNAA}}\) and
\(\text{SpD}_{\scriptscriptstyle \text{noNAA}}\) EMs, exhibited the
lowest global performance among all EMs, performing even worse than the
simplest \(\text{PAN}_{\scriptscriptstyle \text{noNAA}}\) EM
(\autoref{fig:holistic_OM1}). Similar results were found at the regional
scale (Figs S14-S15).

\hypertarget{performance-under-the-low-movement-scenario}{%
\subsubsection{Performance under the Low Movement
Scenario}\label{performance-under-the-low-movement-scenario}}

Short-term and long-term performance showed that
\(\text{EMs}_{\scriptscriptstyle \text{NAA}}\) performed more similarly
to the baseline EM (Figs S26-S27). Specifically,
\(\text{SEP}_{\scriptscriptstyle \text{NAA}}\) and
\(\text{SpD}_{\scriptscriptstyle \text{NAA}}\) EMs exhibited nearly
identical performance to the baseline EM in terms of catch, \(F\), and
\(SSB\) (Figs S26-S27). In contrast,
\(\text{EMs}_{\scriptscriptstyle \text{noNAA}}\), although achieving
higher catch in the short term, led to significant declines in \(SSB\)
in the long term at both regional and global scales, regardless of model
complexity or fishing history (Figs S26-S27). Similar patterns in AACV,
\(\text{P}_{\text{SSB<SSB}_{\text{MSY}}}\),
\(\text{P}_{\text{F>F}_{\text{MSY}}}\), \(\text{OFD}_\text{terminal}\),
and \(\text{OFG}_\text{terminal}\) were observed under the high movement
scenario (Figs S28-S32). These results consistently demonstrate that
\(\text{EMs}_{\scriptscriptstyle \text{NAA}}\) outperformed
\(\text{EMs}_{\scriptscriptstyle \text{noNAA}}\), with the most complex
EMs with NAA generally showing superior performance compared to other
EMs (Figs S33-S35).

\hypertarget{discussion}{%
\section{Discussion}\label{discussion}}

Our findings confirm that management outcomes generally improve when NAA
random effects are included, whereas excluding them from assessment
models leads to undesirable outcomes both regionally and globally.
Earlier claims have suggested that NAA deviations can represent fish
migration in stock assessment models \citep{Gudmundsson2012}. Our study
suggests that the improved management outcomes likely result from the
ability of NAA random effects to implicitly account for movement
variation, in accordance with the understanding that NAA random effects
serve as a `catch-all' to absorb process variation from confounding
processes such as fishing selectivity and natural mortality
\citep{Li2024}.

\hypertarget{improved-parameter-estimation-with-naa-random-effects}{%
\subsection{Improved Parameter Estimation with NAA Random
Effects}\label{improved-parameter-estimation-with-naa-random-effects}}

It is common for one modeled process to alias others in state-space
models \citep{Stock2021b}. Our study demonstrates that NAA random
effects not only alias unmodeled movement dynamics but also mitigate the
consequences of model misspecification caused by the lack of a movement
component in the assessment model. We found that NAA random effects
absorbed movement process variation. Without NAA random effects, mean
recruitment and recruitment variability increased significantly to
compensate for movement. When NAA random effects are included, the
movement-related uncertainty is partitioned between recruitment and NAA.
This partitioning reduces recruitment overestimation and attributes some
of the movement uncertainty to NAA random effects, as evidenced by
decreased recruitment variability and increased NAA variability.
Consequently, \(\text{EMs}_{\scriptscriptstyle \text{NAA}}\) provided
more accurate recruitment estimates, both regionally and globally. For
example, estimates of global recruitment and recruitment in region 2
showed significant improvements in
\(\text{EMs}_{\scriptscriptstyle \text{NAA}}\) compared to models
without them. A similar pattern was observed for \(SSB\), with global
\(SSB\) and region 2 \(SSB\) estimates being more accurate in
\(\text{EMs}_{\scriptscriptstyle \text{NAA}}\) compared to
\(\text{EMs}_{\scriptscriptstyle \text{noNAA}}\). This improvement may
indirectly enhance the accuracy of other parameter estimates, thereby
improving management outcomes.

\hypertarget{improved-management-performance-with-naa-random-effects}{%
\subsection{Improved Management Performance with NAA Random
Effects}\label{improved-management-performance-with-naa-random-effects}}

We showed that including NAA random effects generally improved
management outcomes in nearly every aspect of performance metrics. By
contrast, excluding NAA random effects impeded management outcomes by
producing \(F\) that induced overfishing, partially due to recruitment
overestimation. This overestimation led to higher catch but decreased
\(SSB\), both regionally and globally, in the short term. The absence of
NAA random effects caused more pronounced issues in the long term. For
instance, consistently high \(F\) driven by optimistic estimation of
population size and productivity led to significant reductions in
\(SSB\), while the gain in catch diminished. Compared to the baseline
EM, we found global \(SSB\) declined by up to 20\% and 15\% under high
and low movement scenarios, respectively, without corresponding
increases in catch. Furthermore, mean recruitment and recruitment
variability was amplified in
\(\text{EMs}_{\scriptscriptstyle \text{noNAA}}\), leading to inaccurate
estimation of population size and productivity, thereby impacting the
stability of catch and \(SSB\) in both the short term and long term.
Issues of using \(\text{EMs}_{\scriptscriptstyle \text{noNAA}}\) were
further compounded when assessments were conducted every 3 years, as
errors of population estimates, reference points, and catch advice
accumulated during the interim years, ultimately degrading long-term
management performance.

Under low rates of movement, our results indicate that NAA random
effects have the potential to account for movement and achieve
satisfactory management outcomes. This is echoed by the nearly identical
performance of \(\text{SEP}{\scriptscriptstyle \text{NAA}}\) and
\(\text{SpD}{\scriptscriptstyle \text{NAA}}\) EMs compared to the
baseline EM. However, under high movement rates, NAA random effects may
struggle to fully compensate for the consequences of model
misspecification due to the absence of explicitly modeled movement. For
instance, while \(\text{EMs}_{\scriptscriptstyle \text{NAA}}\) achieved
desirable management outcomes in the short term, in the long term global
\(SSB\) still decreased by 5-10\%, with region 1 experiencing a 0-25\%
reduction and region 2 a 0-30\% reduction, depending on historical
fishing pressure and the spatial complexity of the assessment model. In
addition, \(\text{EMs}_{\scriptscriptstyle \text{NAA}}\) exhibited
limited ability to prevent overfishing, as reflected by even relatively
complex EMs approaching overfishing thresholds in the terminal year.

\hypertarget{benefits-of-correct-spatial-structure}{%
\subsection{Benefits of Correct Spatial
Structure}\label{benefits-of-correct-spatial-structure}}

Our study demonstrates that the benefits of NAA random effects are
maximized when combined with correct spatial structures in stock
assessment models. Accurate spatial boundaries and sufficient data
quantity and quality for spatial disaggregation are critical for
achieving management objectives \citep{Berger2021}. Conversely,
mismatches between population structures and management units can result
in adverse management outcomes
\citep{Cope2011, Hintzen2015, Kerr2017, Berger2021}. Spatial
heterogeneity can arise from spatially heterogeneous biological
characteristics such as recruitment, growth, maturity, and movement
dynamics, or fishery characteristics such as fishing mortality and
selectivity patterns. In our study, such heterogeneity created barriers
for assessment models with incorrect spatial structure to achieve
successful management, as demonstrated by increased risks of localized
depletion due to inappropriate allocation of catch.

Limitations of incorrect spatial structure were exemplified by the poor
performance of the \(\text{PAN}_{\scriptscriptstyle \text{NAA}}\) EM in
our study, even when catch allocation was informed using the proportion
of total survey biomass in each management area, a method regarded as
best practice for maximizing system yield when the underlying spatial
structure is unknown \citep{Bosley2019}. The primary challenge in our
scenario arises from the interaction of natal homing and source-sink
dynamics under high movement rates, which reduces the effectiveness of
this approach. In our high movement scenario, region 1 (source region
with high productivity) disproportionately contributes to global
recruitment but does not retain a large portion of its recruits.
Instead, many recruits and subsequent fish at age 2+ migrate to region 2
(sink region with low productivity) outside the spawning season and
contribute to the fishery there. If catch allocation is based solely on
\(SSB\) or recruitment within each region, this approach will likely
lead to an underestimation of the contribution of the source region to
the fishery in the sink region and overly optimistic estimation about
the productivity from the source region. However, due to the substantial
emigration of recruits to the sink region, the actual biomass available
to the fishery in the source region should be lower than the survey
suggests. This mismatch can result in overharvesting in the source
region and underharvesting in the sink region, ultimately degrading
management outcomes and reducing the likelihood of achieving maximum
sustainable yield, as evidenced by significantly high \(F\) (combined
with risk of overfishing) in region 1 with low \(F\) in region 2
(combined with not overfishing) in our study. In addition, fishing
selectivity patterns differ between regions in our study, while assuming
a `global' fishing selectivity pattern in a PAN EM may lead to
overestimation in one region and underestimation in another. This
mismatch likely perturbs model performance, resulting in incorrect
estimation of fishing status and suboptimal management advice. Even
though management performance may be sustained, PAN EMs should be
applied with caution, as they may lead to localized depletion
\citep{Punt2015, Bosley2019, Bosley2022} and loss of biocomplexity,
ultimately affecting population resilience and stability, particularly
in the face of climate change and human perturbations \citep{Kerr2017}.

The \(\text{FAA}_{\scriptscriptstyle \text{NAA}}\) EM outperformed the
\(\text{PAN}_{\scriptscriptstyle \text{NAA}}\) EM in our study by
providing robust estimates of population status and productivity, with
regional catch advice better reflecting the underlying spatial
heterogeneity in fishing patterns. Compared to the
\(\text{PAN}_{\scriptscriptstyle \text{NAA}}\) EM, the
\(\text{FAA}_{\scriptscriptstyle \text{NAA}}\) EM with spatial structure
modeled implicitly delivered more desirable management outcomes in our
study. We found that the \(\text{FAA}_{\scriptscriptstyle \text{NAA}}\)
EM also outperformed the baseline EM by providing more stable catches at
both regional and global scales, consistent with Punt et al.~(2017).
This highlights their potential as a cost-effective approach for
balancing model complexity and practicality. The fleets-as-areas
approach assumes a single homogenous population while attributing
regional differences in age compositions to fishing selectivity. In some
cases, this flexibility in modeling spatially distinct fleet patterns
helps absorb movement process variation by attributing differences in
age compositions resulting from movement dynamics to fishing selectivity
\citep{Cope2011, Berger2012, O'Boyle2016}. For example, previous studies
have reported dome-shaped selectivity as a result of model behavior that
helps capture spatial heterogeneity in population structure and fishing
with low rates of movement \citep{Waterhouse2014}. Incorporating
time-varying process error into selectivity has been shown to provide
precise estimates of derived and management quantities when spatial
patterns due to age-specific movement are unmodeled \citep{Lee2017}.
This highlights an avenue for future research to explore modeling
selectivity as random effects (in addition to NAA random effects) to
better account for movement dynamics and improve model accuracy.

\(\text{SEP}{\scriptscriptstyle \text{NAA}}\) and
\(\text{SpD}{\scriptscriptstyle \text{NAA}}\) EMs, which explicitly
model spatial structure, demonstrate promising potential in improving
management outcomes. This is evident from their comparable management
performance to the baseline EM in this study. However, their performance
was degraded under the high movement scenario, as they do not fully
capture the movement dynamics. Incorporating spatially referenced
parameters in an assessment model is generally ideal when assumptions of
homogeneity in biology and fisheries are violated, as it accounts for
the spatial structure more effectively. With increasing model complexity
and flexibility, assessment models may better account for the net
effects of movement on observations of population structure at local
scales, effectively addressing spatial structure and movement dynamics.
The primary difference between SEP and SpD EMs lies in how biological
reference points were determined. SpD EMs rely on global SPR-based
reference points, which are weighted by the estimated recruitment within
each region, whereas SEP EMs calculate reference points independently
for each region. In most cases, our results showed minimal differences
in management performance between these two EMs. However, under the
contrasting fishing history, distinct patterns of fishing mortality
impacted recruitment estimates in each region, thereby affecting the
recruitment-based weights that determined the influence of each region
on the global SPR-based reference point. Consequently, the \(F\)
reference points provided by SEP and SpD EMs differed, resulting in
variations in regional management performance, while global management
performance remained largely identical. For instance, source regions
with high recruitment may dominate the global SPR-based reference point,
even if their actual contribution to the fishery is diminished by
emigration. This can lead to inappropriate catch allocations, causing
underharvesting in source regions (due to overestimated recruitment
weights) and overharvesting in sink regions (due to underestimated
recruitment weights), as evidenced in our case where the
\(\text{SpD}_{\scriptscriptstyle \text{NAA}}\) EM resulted in
underfishing for region 1 and overfishing for region 2 when movement
rates were high.

Our results demonstrate that estimating movement in spatial state-space
assessment models by treating movement as random effects with prior
information (i.e.~\(\text{SpE}_{\scriptscriptstyle \text{NAA,Est}}\))
could be a promising solution to better meet management objectives. This
approach showed improved accuracy in parameter estimation and yielded
more favorable management outcomes compared to EMs that excluded
movement. However, assessment models estimating movement often rely on
strong priors, which may be challenging to obtain due to the limited
availability of tagging, genetic, or stable isotope data. While testing
the effects of prior distributions on movement parameter estimation is
beyond the scope of this study, we recommend that future research
investigate the sensitivity of model performance to varying qualities of
prior information in movement estimation.

\hypertarget{caveats-and-future-research-recommendations}{%
\subsection{Caveats and Future Research
Recommendations}\label{caveats-and-future-research-recommendations}}

We note that several assumptions made in our study could influence the
overall findings. For instance, all fish were assumed to be
reproductively isolated and to exhibit natal homing spawning migrations.
This assumption ensured that surveys conducted during the spawning
season had the best opportunity to collect information on \(SSB\) for
each population. However, mismatches between the timing of surveys and
the spawning season could potentially affect management outcomes.
Second, movement rates were assumed to be uniform across all ages. While
ontogenetic movement may be biologically plausible in certain cases
\citep{Waterhouse2014, Hanselman2015}, it remains unclear whether a more
complex structure for NAA random effects (e.g., age-specific movement)
could effectively capture differences in age-based movement dynamics. In
this study, we modeled NAA random effects as IID deviations, which may
have limited ability to account for model mis-specification. More
flexible N NAA configurations, such as AR(1) processes across years or
ages, or two-dimensional AR(1) (2D-AR1) structures across both years and
ages, may enhance the ability of NAA random effects to capture more
complex dynamics. Further investigation is needed to assess the broader
applicability of NAA random effects under such scenarios.

Another key assumption was that selectivity pattern in our estimation
models was constrained to a logistic curve. While this assumption
provided NAA random effects with the best opportunity to account for
spatial heterogeneity caused by movement between regions, it restricted
the ability of selectivity patterns to reflect movement dynamics, which
are often characterized by a dome-shaped selectivity pattern in EMs that
do not account for movement \citep{O'Boyle2016}. Moreover, incorporating
random effects into the selectivity process may offer greater
flexibility to account for the underlying movement dynamics, likely
maximizing the benefits of using selectivity to absorb the movement
process variation, despite a potential increased risk of bias in
population estimates and reference points due to model misspecification
\citep{Fisch2023, Li2024}. Thus, the management outcomes of these more
flexible frameworks within state-space models remain poorly understood
and require further investigation \citep{O'Boyle2016}.

Our approach to biological reference points and catch apportionment may
also influence MSE performance by affecting management outcomes rather
than assessment bias. Apportioning global biological reference points to
regions introduces inherent challenges in comparing single-region and
multi-region model outputs or harvest control rules. Thus, further work
is needed to explore pros and cons of using regionally explicit versus
globally apportioned biological reference points (e.g., separate
vs.~spatially disaggregated approaches), particularly under varying
assumptions of movement dynamics among regions.

Finally, fish biology (stock--recruitment relationship) and fleet
characteristics (e.g., fishing selectivity) were assumed to differ
slightly between populations/regions in the present study (see
\autoref{para_tab}). However, relative impacts of these discrepancies,
along with other biological (e.g., growth rates) and fishery
characteristics, on model performance remain unclear and warrant further
investigation to fully understand the potential of random effects in
management outcomes. Our results are indicative of data-rich stocks with
a known source for connectivity dynamics (natal homing), while
alternative data quality and population dynamic assumptions may improve
or degrade the utility of NAA random effects.

\hypertarget{conclusion}{%
\section{Conclusion}\label{conclusion}}

While claims have been made that NAA random effects can account for
other processes including movement, our study found that assessment
models with NAA random effects outperformed their counterparts without
them. However, NAA random effects could not fully capture movement
dynamics, particularly under high levels of movement. Thus, models with
NAA random effects provide a bridge between conventional panmictic stock
assessments with no explicit movement processes and fully spatial
state--space models that incorporate empirical tagging data to estimate
movement rates. In cases where explicitly modeling movement dynamics is
not feasible, NAA random effects offer a practical intermediate
approach. The results of this study add to the growing body of
literature suggesting that incorporating NAA random effects should be
the default starting point in state--space stock assessments, as their
inclusion leads to more accurate stock assessment and
management-relevant estimates.

\hypertarget{acknowledgements}{%
\section{Acknowledgements}\label{acknowledgements}}

This work was funded by NOAA Fisheries Northeast Fisheries Science
Center (50\%), Northwest Fisheries Science Center (25\%), and Alaska
Fisheries Science Center (25\%). We acknowledge the support provided by
Saltwater Inc.~for facilitating Chengxue Li's contribution to this
project. We are grateful to Christopher Legault for conducting NOAA
internal review of the manuscript and to Dana Hanselman for his
insightful contributions to the conceptualization of this study design.
We also thank Microsoft Azure for providing high-performance computing
resources for this study. We would also like to thank the Quahog
Republic and the Midway Trap and Skeet Club for keeping the wheels
adequately greased during the study design development.

\hypertarget{competing-interests-statement}{%
\section{Competing interests
statement}\label{competing-interests-statement}}

One co-author, Daniel R. Goethel, serves as an Associate Editor for
CJFAS, and another co-author, Timothy J. Miller, serves as a Guest
Editor for CJFAS for this special issue.

\hypertarget{credit-authorship-contribution-statement}{%
\section{CRediT authorship contribution
statement}\label{credit-authorship-contribution-statement}}

\textbf{Chengxue Li:} Conceptualization, Methodology, Software, Writing
- original draft, Formal analysis, Visualization.

\textbf{Jonathan J. Deroba:} Conceptualization, Funding acquisition,
Supervision, Writing - review \& editing.

\textbf{Aaron M. Berger:} Conceptualization, Funding acquisition,
Writing - review \& editing.

\textbf{Daniel R. Goethel:} Conceptualization, Funding acquisition,
Writing - review \& editing.

\textbf{Brian J. Langseth:} Conceptualization, Funding acquisition,
Writing - review \& editing.

\textbf{Amy M. Schueller:} Conceptualization, Writing - review \&
editing.

\textbf{Timothy J. Miller:} Software, Writing - review \& editing.

\hypertarget{funding-statement}{%
\section{Funding statement}\label{funding-statement}}

This work was funded by NOAA Fisheries Northeast Fisheries Science
Center (50\%), Northwest Fisheries Science Center (25\%), and Alaska
Fisheries Science Center (25\%).

\hypertarget{data-availability-statement}{%
\section{Data availability
statement}\label{data-availability-statement}}

No empirical data were collected for this study. Data used for this
study were produced through simulation. Code for all simulations is
available at \url{https://github.com/lichengxue/MSE_NAA_Project}.

\hypertarget{tables}{%
\section{Tables}\label{tables}}

\begin{center}
\begin{longtable}{p{2cm} p{10cm} p{3cm}}
\caption{Biological and fishery parameter values used in the operating models ($\text{OM}_{\text{high}}$ and $\text{OM}_{\text{low}}$). When a single value is provided, it applies to both populations/regions. When two values are provided, the first corresponds to either population 1 (bold) or region 1, and the second value corresponds to either population 2 (bold) or region 2. Settings for the closed-loop management strategy evaluation are also provided.} \\
\label{para_tab} \\
\hline
\multicolumn{1}{c}{\textbf{Parameters}} & \multicolumn{1}{c}{\textbf{Description}} & \multicolumn{1}{c}{\textbf{Value}} \tabularnewline
\hline

\multicolumn{3}{l}{\textbf{Biological information}} \tabularnewline
$\alpha_{SR}$ & Productivity parameter & \textbf{7, 5.5} \tabularnewline
$\beta_{SR}$ & Density-dependent recruitment parameter & $1 \times 10^{-4}$ \tabularnewline
$\sigma_R$ & Standard deviation of recruitment random effects (age 1) & 0.5 \tabularnewline
$\sigma_{NAA}$ & Standard deviation of numbers-at-age (age 2+) random effects & 0.2 \tabularnewline
$a_{max}$ & Maximum age (yr) & 12 \tabularnewline
$M$ & Instantaneous natural mortality & 0.2 \tabularnewline
$a_0$ & Age at length = 0 & 0 \tabularnewline
$L_\infty$ & Maximum length (cm) & 90 \tabularnewline
$k$ & Growth rate (yr$^{-1}$) & 0.13 \tabularnewline
$b_1$ & L – W scalar (kg cm$^{-1}$) & $3 \times 10^{-6}$ \tabularnewline
$b_2$ & L – W exponent & 3 \tabularnewline
$m_{50}$ & Age at 50\% maturity & 3.5 \tabularnewline
$m_{slope}$ & Slope of maturity & 1 \tabularnewline

\multicolumn{3}{l}{\textbf{Spatial structure and movement dynamics}} \tabularnewline
$n_{R}$ & Number of regions & 2 \tabularnewline
$n_{season}$ & Number of seasons & 4 \tabularnewline
$f_{SSB}$ & Fraction of the year that passes before spawning occurs & 0.625 \tabularnewline
$\mu_1$ & Movement rate outside the spawning and spawning migration seasons & {0.3, 0.1} \\ & & ($\text{OM}_{\text{high}}$); \\ & & \textbf{0.1, 0.0} \\ & & \textbf{($\text{OM}_{\text{low}}$)} \tabularnewline
$\mu_2$ & Movement rate during the spawning migration season & 1 \tabularnewline
$\mu_3$ & Movement rate during the spawning season & 0 \tabularnewline

$\sigma_\mu$ & Standard deviation of movement rate & 0.1 \tabularnewline
$\rho_\mu$ & Autocorrelation in movement & 0.5 \tabularnewline

\multicolumn{3}{l}{\textbf{Fishery information}} \tabularnewline
$n_{F}$ & Number of fleets in each region & 1 \tabularnewline
$\gamma_{50}$ & Age at 50\% selectivity in fishery & 5.5, 5 \tabularnewline
$\gamma_{slope} $ & Slope of logistic selectivity for fishery & 1 \tabularnewline
$ESS_{C}$ & Effective sample size for fleet catch & 100 \tabularnewline
$CV_{C}$ & Coefficient of variation for fleet catch & 0.1 \tabularnewline
$LL_{C}$ & Likelihood for age compositional data & Multinomial \tabularnewline

\multicolumn{3}{l}{\textbf{Survey information}} \tabularnewline
$n_{I}$ & Number of indices in each region & 1 \tabularnewline
$\zeta_{50}$ & Age at 50\% selectivity in survey & 2 \tabularnewline
$\zeta_{slope}$ & Slope of logistic selectivity for survey & 1 \tabularnewline
$ESS_{I}$ & Effective sample size for survey index & 100 \tabularnewline
$CV_{I}$ & Coefficient of variation for survey index & 0.1 \tabularnewline
$q$ & Survey catchability & 0.2 \tabularnewline
$f_{I}$ & Proportion of the year elapsed when index is observing the population & 0.625 \tabularnewline
$LL_{I}$ & Likelihood for age compositional data & Multinomial \tabularnewline

\multicolumn{3}{l}{\textbf{MSE closed-loop}} \tabularnewline
$\tau_{hist.}$ & Number of years in the historical period & 30 \tabularnewline
$\tau_{feedback}$ & Number of years in the feedback period & 30 \tabularnewline
$\Delta t$ & Assessment interval (yr) & 3 \tabularnewline
$A_{total}$ & Number of total assessments & 10 \tabularnewline
$\Phi$ & A fraction of $F_{40\%}$ used as the harvest control rule & 100\% \tabularnewline

\hline
\end{longtable}
\end{center}


\begin{table}[H]
\centering
\caption{Description of fishing scenarios in the historical period used for operating models ($\text{OM}_{\text{high}}$ and $\text{OM}_{\text{low}}$). The "Adaptive" scenario ($\text{F}_\text{A}$) represents fishing mortality rates for both regions transitioning from high to sustainable levels (2.5\(F_{MSY}\)–\(F_{MSY}\)). The "Unsustainable" scenario ($\text{F}_\text{U}$) represents constantly high fishing mortality rates (2.5\(F_{MSY}\)), and the "Contrasting" scenario ($\text{F}_\text{C}$) represents contrasting fishing mortality rates in each region, with one region experiencing sustainable fishing (\(F_{MSY}\)) and the other unsustainable (2.5\(F_{MSY}\)).
}
\label{F_scenario_table}
\begin{tabular}{p{5cm} p{5cm} p{5cm}}
\toprule
\textbf{Fishing scenarios} & \textbf{Region 1} & \textbf{Region 2} \\
\midrule
Adaptive ($\text{F}_\text{A}$) & H – L (2.5\(F_{MSY}\) – \(F_{MSY}\)) & H – L (2.5\(F_{MSY}\) – \(F_{MSY}\)) \\
Unsustainable ($\text{F}_\text{U}$) & H (2.5\(F_{MSY}\)) & H (2.5\(F_{MSY}\)) \\
Contrasting ($\text{F}_\text{C}$) & H – L (2.5\(F_{MSY}\) – \(F_{MSY}\)) & L – H (\(F_{MSY}\) – 2.5\(F_{MSY}\)) \\
\bottomrule
\end{tabular}
\end{table}


\begin{landscape}
\begin{longtable}{p{3cm} p{6cm} p{2cm} p{3.5cm} p{5.5cm}}  
  \caption{Descriptions of 10 estimation models} \label{EM_tabel} \\
  \toprule
  \textbf{Model} & \textbf{Type} & \textbf{Move} & \textbf{Random effects} & \textbf{Description} \\
  \midrule
  \endfirsthead
  \caption[]{(continued)} \\
  \toprule
  \textbf{Model} & \textbf{Type} & \textbf{Move} & \textbf{Random effects} & \textbf{Description} \\
  \midrule
  \endhead
  \bottomrule
  \endfoot

  PAN$_{\text{NAA}}$ & Panmictic (catch aggregated) & No & NAA, Rec & Catch data aggregated across regions \\
  PAN$_{\text{noNAA}}$ & Panmictic (catch aggregated) & No & \textbf{Rec only} & Catch data aggregated across regions \\
  FAA$_{\text{NAA}}$ & Spatially implicit (fleets-as-areas) & No & NAA, Rec & Multiple fleets account for spatial difference in fleet structure \\
  FAA$_{\text{noNAA}}$ & Spatially implicit (fleets-as-areas) & No & \textbf{Rec only} & Multiple fleets account for spatial difference in fleet structure \\
  SEP$_{\text{NAA}}$ & Separate & No & NAA, Rec & Separate single stock assessment model \\
  SEP$_{\text{noNAA}}$ & Separate & No & \textbf{Rec only} & Separate single stock assessment model \\
  SpD$_{\text{NAA}}$ & Spatially disaggregated & No & NAA, Rec & Spatially disaggregated without movement \\
  SpD$_{\text{noNAA}}$ & Spatially disaggregated & No & \textbf{Rec only} & Spatially disaggregated without movement \\
  SpE$_{\text{NAA,Est}}$ & Spatially explicit & Yes (Est) & NAA, Rec, Move & Spatially explicit with movement estimated using a prior \\
  SpE$_{\text{NAA,Fix}}$ & Spatially explicit & Yes (Fix) & NAA, Rec & Spatially explicit with movement fixed as known \\
  
\end{longtable}
\end{landscape}


\hypertarget{figures}{%
\section{Figures}\label{figures}}

\begin{figure}[H]
\centering
\includegraphics[width=\textwidth]{Movement_dynamics.PNG}
\caption{Movement dynamics for high ($\text{OM}_{\text{high}}$) and low ($\text{OM}_{\text{low}}$) movement scenarios. Blue represents the population originating from the upper region and red represents the population originating from the lower region. The dashed arrow represents movement outside the spawning season, with thickness indicating the magnitude of the movement rate, while the solid curved arrow indicates natal homing movement.}
\label{fig:movement_dynamics}
\end{figure}

\begin{figure}[H]
\centering
\includegraphics[width=\textwidth]{Combined_Figures.PNG}
\caption{Overview of the 10 estimation models (EMs) used in the present study. The description of the model structure is shown in Table~2. The dashed black arrow in the $\text{SpE}_{\scriptscriptstyle \text{NAA,Est}}$ EM indicates that movement outside the spawning season was estimated using a prior, while the solid black arrow in the $\text{SpE}_{\scriptscriptstyle \text{NAA,Fix}}$ EM indicates that movement outside the spawning season was fixed at the true value. The blue curved arrow in $\text{SpE}_{\scriptscriptstyle \text{NAA,Est}}$ and $\text{SpE}_{\scriptscriptstyle \text{NAA,Fix}}$ EMs indicates natal homing movement.}
\label{fig:EM_structure}
\end{figure}

\begin{figure}[H]
\centering
\includegraphics[width=\textwidth]{Performance_first_5_years_relative_OM_1.PNG}
\caption{Relative difference in catch, $F$, and $SSB$ between the baseline EM ($\text{SpE}_{\scriptscriptstyle \text{NAA,Fix}}$) and other EMs over the first 5 years of the feedback period under the high movement scenario ($\text{OM}_{\text{high}}$). Each point represents the median across the 100 realizations, and error bars represent the interquartile range (25th–75th percentiles). EMs with a grey background indicate that NAA random effects were included, while EMs with a white background indicate that NAA random effects were not included. Three fishing scenarios are included: Adaptive ($\mathrm{F_A}$), Unsustainable ($\mathrm{F_U}$), and Contrasting ($\mathrm{F_C}$).}
\label{fig:first_5yr_results}
\end{figure}

\begin{figure}[H]
\centering
\includegraphics[width=\textwidth]{Performance_last_5_years_relative_OM_1.PNG}
\caption{Relative difference in catch, $F$, and $SSB$ between the baseline EM ($\text{SpE}_{\scriptscriptstyle \text{NAA,Fix}}$) and other EMs over the last 5 years of the feedback period under the high movement scenario ($\text{OM}_{\text{high}}$). Each point represents the median across the 100 realizations, and error bars represent the interquartile range (25th–75th percentiles). EMs with a grey background indicate that NAA random effects were included, while EMs with a white background indicate that NAA random effects were not included. Three fishing scenarios are included: Adaptive ($\mathrm{F_A}$), Unsustainable ($\mathrm{F_U}$), and Contrasting ($\mathrm{F_C}$).}
\label{fig:last_5yr_results}
\end{figure}

\begin{figure}[H]
\centering
\includegraphics[width=\textwidth]{Catch_variation_OM1.PNG}
\caption{Average annual catch variation (AACV) of each EM at both regional and global scales under the high movement scenario ($\text{OM}_{\text{high}}$). Each point represents the median across the 100 realizations (one value per realization), and error bars represent the interquartile range (25th–75th percentiles). EMs with a grey background indicate that NAA random effects were included, while EMs with a white background indicate that NAA random effects were not included. Three fishing scenarios are included: Adaptive ($\mathrm{F_A}$), Unsustainable ($\mathrm{F_U}$), and Contrasting ($\mathrm{F_C}$).}
\label{fig:aacv}
\end{figure}

\begin{figure}[H]
\centering
\includegraphics[width=\textwidth]{Probability_Overfished_Overfishing_OM1_MSY.PNG}
\caption{Probability of $SSB < SSB_{MSY}$ and $F > F_{MSY}$ for each EM over the 30-year feedback period under the high movement scenario ($\text{OM}_{\text{high}}$). Each point represents the median across the 100 realizations (one value per realization), and error bars represent the interquartile range (25th–75th percentiles). EMs with a grey background indicate that NAA random effects were included, while EMs with a white background indicate that NAA random effects were not included. Three fishing scenarios are included: Adaptive ($\mathrm{F_A}$), Unsustainable ($\mathrm{F_U}$), and Contrasting ($\mathrm{F_C}$).}
\label{fig:Overfished_OM1}
\end{figure}

\begin{figure}[H]
\centering
\includegraphics[width=\textwidth,height=\textheight]{Holistic_OM_1.PNG}
\caption{Overall performance of each EM at the global scale under the high movement scenario ($\text{OM}_{\text{high}}$), including the following relative performance metrics: 1) long-term catch ($\text{Catch}_{\text{long}}$); 2) short-term catch ($\text{Catch}_{\text{short}}$); 3) long-term $SSB$ ($SSB_{\text{long}}$); 4) short-term $SSB$ ($SSB_{\text{short}}$); 5) probability of $F < F_{MSY}$ ($\text{P}_{F<F_{MSY}}$); 6) probability of $SSB > SSB_{MSY}$ ($\text{P}_{SSB>SSB_{MSY}}$); 7) overfishing status in the terminal year ($\text{OFG}_{\text{terminal}}$); 8) overfished status in the terminal year ($\text{OFD}_{\text{terminal}}$); and 9) average annual catch variation (AACV). Bars represent the median score across the 100 realizations, and error bars represent the interquartile range (25th–75th percentiles). The number displayed in each panel represents the total score (sum of all medians for that EM) for the corresponding EM under that fishing scenario, with higher totals indicating better overall performance. The black dashed line separates the performance of each EM. EMs with a grey background indicate that NAA random effects were included, while EMs with a white background indicate that NAA random effects were not included. Three fishing scenarios are included: Adaptive ($\mathrm{F_A}$), Unsustainable ($\mathrm{F_U}$), and Contrasting ($\mathrm{F_C}$).}
\label{fig:holistic_OM1}
\end{figure}

\pagebreak

\bibliography{paper}

\pagebreak

\end{document}
